% !TEX program = xelatex
% =====================================================================
% a9lim - resume
% Conventional one-page research resume. Build with ./resume/build.sh.
% =====================================================================

\documentclass[10pt,letterpaper]{article}

\usepackage[letterpaper,top=0.6in,bottom=0.6in,left=0.6in,right=0.6in]{geometry}
\usepackage{fontspec}
\usepackage{xcolor}
\usepackage{titlesec}
\usepackage{enumitem}
\usepackage{hyperref}
\usepackage{tabularx}
\usepackage{microtype}

\setmainfont{Recursive}[
  Path           = fonts/,
  Extension      = .ttf,
  UprightFont    = *-Regular,
  BoldFont       = *-Bold,
  ItalicFont     = *-Italic,
  BoldItalicFont = *-BoldItalic,
  Ligatures      = TeX,
]

\definecolor{water}{HTML}{000080}

\hypersetup{
  colorlinks  = true,
  urlcolor    = water,
  linkcolor   = water,
  pdftitle    = {a9lim's Resume},
  pdfauthor   = {a9lim},
  pdfsubject  = {Resume},
  pdfkeywords = {mechanistic interpretability, AI alignment, sparse crosscoders, Abstract-CoT, Python, PyTorch, WebGPU, WGSL},
}

\pagestyle{empty}
\setlength{\parindent}{0pt}
\setlength{\parskip}{0pt}
\setlength{\tabcolsep}{0pt}
\raggedbottom

% Compact, conventional section headings with a single restrained rule.
\titleformat{\section}
  {\large\bfseries}
  {}{0pt}{}
\titlespacing*{\section}{0pt}{2.5pt}{2.5pt}

% Standard resume bullets: readable hanging indent, no decorative markers.
\setlist[itemize]{
  leftmargin=10pt,
  label=\raisebox{0.2ex}{\scriptsize\textbullet},
  labelsep=2.5pt,
  itemsep=1pt,
  topsep=1pt,
  parsep=0pt,
  partopsep=0pt
}

\newcommand{\entry}[4]{%
  \noindent\begin{tabularx}{\linewidth}{@{}X r@{}}
    \textbf{\href{#2}{#1}} & \small #4 \\
    \textit{#3} & \\
  \end{tabularx}%
}

\newcommand{\plainentry}[3]{%
  \noindent\begin{tabularx}{\linewidth}{@{}X r@{}}
    \textbf{#1} & \small #3 \\
    \textit{#2} & \\
  \end{tabularx}%
}

\begin{document}

% Header
\noindent
\begin{minipage}[b]{0.24\linewidth}
  {\fontsize{23}{24}\selectfont\bfseries a9lim}
\end{minipage}%
\hfill
\begin{minipage}[b]{0.74\linewidth}
  \raggedleft\small
  Ann Arbor \enspace\textbar\enspace
    \href{mailto:mx@a9l.im}{mx@a9l.im} \enspace\textbar\enspace
    \href{https://a9l.im}{a9l.im} \enspace\textbar\enspace
    \href{https://github.com/a9lim}{github.com/a9lim}
\end{minipage}
\par\vspace{2.5pt}
{\small
Independent interpretability and alignment researcher, and open-source developer.
\par}

\section*{Education}

\plainentry{Singapore American School}{High School Diploma - Weighted GPA 4.5}{Aug 2019 - May 2023}

\plainentry{University of California, San Diego}{Bachelor of Science in Mathematics - GPA 3.73 - Major GPA 3.86}{Sep 2023 - Mar 2026}

\section*{Research}

\entry{Introspection via Kaomoji}{https://a9l.im/blog/introspection-via-kaomoji}{Study on legibility of emotional representation in LLMs - Python, HF Transformers}{Apr 2026 - Jun 2026}
\begin{itemize}
  \item Investigated correspondence between kaomoji used by LLMs and functional emotional states.
  \item Found that the kaomoji could predict the emotional content of the prompt with 80.6\% accuracy on Gemma-4-31b-it.
  \item Aligned hidden states across model families to analyze shared geometry.
  \item Built data collection tooling and published the resulting corpus.
\end{itemize}

\entry{Elapsed-Time Encoding in LLMs}{https://github.com/a9lim/time-experiment}{Ongoing study on representations of elapsed time within LLMs - Python, HF Transformers}{Jun 2026 - Present}
\begin{itemize}
  \item Used linear probes and logits to investigate the representation of elapsed time within LLMs over a conversation without timestamps.
  \item Showed elapsed time is tracked within LLMs and linearly proportional to context length using probes.
  \item Replicated findings of Discrete Minds in a Continuous World (Wang et al.).
\end{itemize}

\entry{Neuralese as a Conlang}{https://github.com/a9lim/acot-experiment}{Ongoing study on the legibility of latent reasoning - Python, PyTorch}{Jul 2026 - Present}
\begin{itemize}
  \item Train a model to reason via Abstract Chain-of-Thought (Ramji et al.) to analyze the structure and grammar of abstract tokens.
  \item A preliminary run on a small model has so far not successfully trained it to use abstract tokens to reason.
\end{itemize}

\entry{Block Sparse Crosscoders}{https://github.com/a9lim/block-crosscoder-experiment}{Ongoing study on cross-layer multidimensional features in LLMs - Python, PyTorch}{Jul 2026 - Present}
\begin{itemize}
  \item Synthesize recent work on Block Sparse Featurizers (Fel et al.) and Sparse Crosscoders (Lindsey et al.) to attempt extraction of cross-layer multidimensional features.
\end{itemize}

\section*{Projects}

\entry{Saklas}{https://github.com/a9lim/saklas}{Mechanistic interpretability workbench - Python, HF Transformers, PyTorch}{Apr 2026 - Present}
\begin{itemize}
  \item Built a local interpretability workbench compatible with most open-weight LLMs and both CUDA and MPS.
  \item Implemented curved manifold extraction from sets of concepts.
  \item Implemented Jacobian lens readouts and steering, and nonstandard roles in chat templates.
  \item Published a Python API, dashboard, and OpenAI/Ollama-compatible server.
\end{itemize}

\entry{Geon}{https://a9l.im/geon}{Relativistic N-body simulation - JavaScript, WebGPU}{Feb 2026 - Present}
\begin{itemize}
  \item Built a WebGPU simulation that models N-body physics with charged rotating massive particles.
  \item Implemented finite-speed force propagation and Einstein–Infeld–Hoffmann dynamics.
  \item Implemented force-coupled scalar fields with self-gravity.
\end{itemize}

\plainentry{Other Projects}{Additional projects and tooling I've worked on - Rust, Python, JavaScript, WebGPU}{Feb 2026 - Present}
\begin{itemize}
  \item \textbf{\href{https://github.com/a9lim/ogdoad}{Ogdoad}:} Built a library to explore Clifford algebras over the
  surreal numbers and Conway's combinatorial games and investigate open problems.
  \item \textbf{\href{https://github.com/a9lim/gaslamp}{Gaslamp}:} Built a
  cross-agent CLI for Claude Code and Codex to dispatch each other as subagents.
  \item \textbf{\href{https://a9l.im/sims}{Interactive educational simulations}:} Built seven more browser simulations across physics,
  finance, biology, political science, and comparative religion.
\end{itemize}

\section*{Technical Skills}

{\small
\textbf{Research:} Experiment design and execution, agent orchestration.\\
\textbf{ML:} Python, HF Transformers, PyTorch, NumPy/SciPy.\\
\textbf{Programming:} Java, Rust, JavaScript, WebGPU, Git.
}

\end{document}
